\documentclass[paper=screen,mode=present,style=zysimple]{powerdot}

\usepackage{fancyvrb,verbatim,listings}
\usepackage{proof}
\usepackage[all]{xypic,xy}

\usepackage{latexsym,stmaryrd,amsmath,amsfonts}

%% \newtheorem{definition}{Definition}
%% \newtheorem{theorem}{Theorem}
%% \newtheorem{lemma}{Lemma}
\newtheorem{prop}{Proposition}
%% \newtheorem{corollary}{Corollary}
%% \newenvironment{proof}{{\sc Proof . }}{{\hfill $\square$}}
%% \newcommand{\myqed}{\hfill $\square$}
%% \newcommand{\ped}{\vspace{-1.6em}}
%% \newcommand{\note}[1]{\noindent\fbox{\parbox{\textwidth}{\em #1}}}

\newcommand{\abind}[2]{\mathtt{bind} \ (\, {#1} \,)  \ \mathtt{in} \ {#2}}
\newcommand{\acase}{\mathtt{case}}
\newcommand{\adv}{\mathcal{A}}
\newcommand{\aspecto}[1]{\xrightarrow{\; {#1} \;}}
\newcommand{\asplinear}{\mathfrak{l}}
\newcommand{\aspmodal}{\mathfrak{m}}
\newcommand{\aspnormal}{\mathfrak{n}}
%
\newcommand{\betareduce}{\leadsto_\beta}
\newcommand{\betareduceto}{\twoheadrightarrow_\beta}
\newcommand{\bit}{\mathtt{B}}
\newcommand{\bitone}{\mathtt{B}_1}
\newcommand{\bitzero}{\mathtt{B}_0}
%
\newcommand{\caseto}{\triangleright}
\newcommand{\cbind}[3]{\mathtt{bind} \ {#1} = {#2} \ \mathtt{in} \ {#3}}
\newcommand{\ccase}[4]{\mathtt{case} \; {#1} \; \mathtt{of} \; \cnil \caseto {#2}; \;
                                          \bitzero x_0 \caseto {#3}; \; \bitone x_1 \caseto {#4}}
\newcommand{\cequiv}{\simeq}
\newcommand{\cif}[3]{\mathtt{if} \; {#1} \; \mathtt{then} \; {#2} \; \mathtt{else} \; {#3}}
\newcommand{\ctrue}{\mathtt{true}}
\newcommand{\cfalse}{\mathtt{false}}
\newcommand{\cfix}{\mathtt{fix}\,}
\newcommand{\cfixin}[2]{\mathtt{fix}\; {#1}\; \mathtt{in} \; {#2}}
\newcommand{\cfold}{\mathtt{fold}\,}
\newcommand{\chd}{\mathtt{hd}}
\newcommand{\cinj}{\mathtt{inj}}
\newcommand{\clift}[1]{\lfloor{#1}\rfloor}
\newcommand{\clet}[3]{\mathtt{let} \ {#1} = {#2} \ \mathtt{in} \ {#3}}
\newcommand{\cletst}[3]{\mathtt{liftlet} \ {#1} = {#2} \ \mathtt{in} \ {#3}}
\newcommand{\cletrec}[3]{\mathtt{letrec} \ {#1} = {#2} \ \mathtt{in} \ {#3}}
\newcommand{\cnil}{\mathtt{nil}}
\newcommand{\conc}{\! \bullet \!}
\newcommand{\concat}{\pmb{\mathit{conc}}}
\newcommand{\context}{\mathbb{C}}
\newcommand{\cop}{\;\mathtt{op}\;}
\newcommand{\cpair}[1]{\langle{#1}\rangle}
\newcommand{\cproj}{\mathtt{proj}}
\newcommand{\crand}{\mathtt{rand}}
\newcommand{\crec}{\mathtt{rec}}
\newcommand{\croll}{\mathtt{roll}\,}
\newcommand{\crule}[1]{\mbox{{\textsl{#1}}}}
\newcommand{\cslrtype}{\Vdash}
\newcommand{\csucc}{\mathtt{S}}
\newcommand{\ctl}{\mathtt{tl}}
\newcommand{\cut}{\pmb{\mathit{cut}}}
\newcommand{\cval}{\mathtt{return}}
\newcommand{\cwhile}[2]{\mathtt{while}\; {#1} \; \mathtt{do} \; {#2} \; \mathtt{od}}
\newcommand{\cvalof}{\pmb{valueof}}
\newcommand{\cunfold}{\mathtt{unfold}\,}
\newcommand{\cunit}{\mathtt{unit}}
\newcommand{\cunroll}{\mathtt{unroll}\,}
\newcommand{\cupdate}{\pmb{update}}
\newcommand{\czero}{\mathtt{O}}
%
\newcommand{\deduces}{\vdash}
\newcommand{\defeq}{\stackrel{\mathrm{def}}{=}}
\newcommand{\delay}[1]{\left\lfloor{#1}\right\rfloor}
\newcommand{\distr}{\mathcal{D}}
%
\newcommand{\eagerreduce}{\xrightarrow{\mathrm{eager}}}
\newcommand{\Eagerreduce}{\stackrel{\mathrm{eager}}{\Reduce}}
\newcommand{\eqtest}{\stackrel{\mbox{\tiny ?}}{=}}
\newcommand{\Eval}[1]{{\left\llbracket{#1}\right\rrbracket}}
\newcommand{\evalcontext}{\mathbb{E}}
\newcommand{\Example}{\emgreen{\pmb{\textsc{Example}\;}}}
\newcommand{\expansion}[1]{\tilde{{#1}}}
%
\newcommand{\ford}{\mathrm{ord}}
%
\newcommand{\getval}{\leftarrow}
%
\newcommand{\head}{\pmb{\mathit{hd}}}
\newcommand{\termhole}{[\;]}
\newcommand{\hoaretriple}[3]{\{{#1}\} \; {#2} \; \{{#3}\}}
%
\newcommand{\imply}{\Rightarrow}
\newcommand{\indist}{\simeq}
\newcommand{\iteration}{\pmb{\mathit{it}}}
%
\newcommand{\lazyreduce}{\xrightarrow{\mathrm{lazy}}}
\newcommand{\Lazyreduce}{\lazyreduce\!\!\!\!\!\to}
\newcommand{\ldot}{\, . \, }
\newcommand{\linearto}{\multimap}
\newcommand{\logrel}{(\mathcal{R}_\tau)_{\mbox{$\tau$ type}}}
\newcommand{\lr}[1]{\mathcal{R}_{#1}}
\newcommand{\lub}{\sqcup}
%
\newcommand{\metalambda}{\underline{\lambda}}
\newcommand{\modal}{\square}
%
\newcommand{\nat}{\mathbb{N}}
\newcommand{\nReduce}{\nrightarrow\!\!\!\!\to}
\newcommand{\nreduceto}{\nrightsquigarrow}
\newcommand{\nullenv}{\varnothing}
%
\newcommand{\orc}{\mathcal{O}}
\newcommand{\oslrtype}{\Vdash}
%
\newcommand{\porder}{\sqsubseteq}
\newcommand{\prefix}{\pmb{\mathit{pref}}}
\newcommand{\proba}{\mathbf{Pr}}
%
\newcommand{\randlen}{\pmb{\mathit{rl}}}
\newcommand{\randnum}{\pmb{\mathit{rn}}}
\newcommand{\randstr}{\pmb{\mathit{rs}}}
\newcommand{\reduce}{\leadsto}
\newcommand{\reduceto}{\rightsquigarrow\hspace*{-6px}\shortrightarrow}
\newcommand{\Reduce}{\to\!\!\!\to}
\newcommand{\Reduceto}{\Downarrow}
%
\newcommand{\set}[1]{\{{#1}\}}
\newcommand{\setbits}{\mathbb{B}}
\newcommand{\slrtype}{\vdash}
\newcommand{\splt}{\pmb{\mathit{split}}}
\newcommand{\suffix}{\pmb{\mathit{suff}}}
\newcommand{\sysname}{\pmb{\mathcal{C}}}
%
\newcommand{\tail}{\pmb{\mathit{tl}}}
\newcommand{\tbits}{\mathsf{Bits}}
\newcommand{\tbool}{\mathsf{Bool}}
\newcommand{\tnat}{\mathsf{Nat}}
\newcommand{\tcomp}{\mathsf{T}}
\newcommand{\tlift}[1]{\lfloor{#1}\rfloor}
\newcommand{\tlist}{\mathsf{List}}
\newcommand{\tself}{\mathsf{self}}
\newcommand{\tstore}{\mathsf{Store}}
\newcommand{\tstream}{\mathsf{Stream}}
\newcommand{\tloc}{\mathsf{Loc}}
\newcommand{\ttcase}{\mathtt{case}}
\newcommand{\ttif}{\mathtt{if}}
\newcommand{\ttthen}{\mathtt{then}}
\newcommand{\ttelse}{\mathtt{else}}
\newcommand{\ttin}{\mathtt{in}}
\newcommand{\ttlet}{\mathtt{let}}
\newcommand{\ttof}{\mathtt{of}}
\newcommand{\tunit}{\mathsf{Unit}}
\newcommand{\Type}{\mathcal{T}\!\mathit{ype}}
%
\newcommand{\vbool}{\pmb{\mathbb{B}}}
\newcommand{\vff}{\mathfrak{f}\!\mathfrak{f}}
\newcommand{\vnat}{\pmb{\mathrm{N}}}
\newcommand{\vnil}{\epsilon}
\newcommand{\vtt}{\mathfrak{t}\!\mathfrak{t}}
%
\newcommand{\weakpre}{\pmb{wp}}




\pdsetup{lf={Theorey and Practice of Functional Programming}}

\title{{\small \textsc{Theory and Practice of Functional Programming}}
\\[12pt]
{\Large \color{TitleColor} Lambda-Calculus}
}
\author{\small \textbf{Dr. ZHANG Yu}
  \\[2pt] \small Institute of Software, Chinese Academy of Sciences 
  \\[5pt] \includegraphics[height=.04\slideheight]{iscas.eps} 
  \qquad \qquad 
  \includegraphics[height=.04\slideheight]{cas.eps}
}
\date{\scriptsize GUCAS, Beijing}

\newcommand{\RED}{\mathsf{RED}}

\begin{document}

\maketitle 

%% \section{}

\begin{slide}{The untyped $\lambda$-calculus}
$\lambda$-calculus is the fundamental model for programming languages. 
It itself is seen as a language that is built on very simple syntax, yet has powerful computing capability. 
\begin{itemize}
\item In $\lambda$-calculus, the basic units are \emblue{expressions} or \emblue{terms}.
\begin{wrap}\vspace*{-0.5em}
\[
\emblue{e, e', \ldots :: = x \mid \lambda x \ldot e \mid e \, e'}
\vspace*{-1em}
\]
\end{wrap}
where $x$ ranges over an infinite set of variables.
\end{itemize}
\end{slide}

\begin{slide}{The untyped $\lambda$-calculus}
\begin{itemize}
\item \emred{Abstraction}: $\lambda x . e$ --- a function where $x$ is the argument and $e$ is the function body.
\begin{itemize}
\item The variable $x$ is \emblue{bound} in the expression.  
\item \Example ($\lambda$-abstractions) 
\[
\begin{array}{l@{\quad}l}
\lambda x \ldot x + 1 & \mbox{the increment function} \\[0.2em]
\lambda x . \lambda y \ldot x + y & \mbox{the addition} \\[0.2em]
\lambda f . \lambda g . \lambda x \ldot f (g \, x) & \mbox{function composition} 
\end{array}
\]
\empurple{We shall see how to program ``+'' and numerals later.}
\item The scope of $\lambda$ extends as far as possible to the right, e.g., $\lambda x \ldot e \, e'$ is 
actually $\lambda x \ldot (e \, e')$. 
\item We also  write $\lambda x \ldot \lambda y \ldot e$ as $\lambda x y \ldot e$.
\end{itemize}
\end{itemize}
\end{slide}

\begin{slide}{The untyped $\lambda$-calculus}
\begin{itemize}
\item \emred{Application}: $e \, e'$ --- compute function $e$ with the (concrete) argument $e'$, 
  or apply function $e$ to the argument $e'$.
\[
\begin{array}{l}
(\lambda x \ldot x + 1) \, 2 \\[0.2em]
(\lambda x . \lambda y \ldot x + y) \, 1 \\[0.2em]
(\lambda f . \lambda g . \lambda x \ldot f (g \, x)) \, (\lambda x \ldot x + 1) \, (\lambda x \ldot x * 2)
\end{array}
\]
\begin{itemize}
\item Applications are left-associative, i.e., $e_1 \, e_2 \, e_3$ stands for $(e_1 \, e_2) \, e_3$.
\end{itemize}
\end{itemize}
\end{slide}

\begin{slide}[method=direct]{Free and bound variables}
\begin{itemize}
\item \emblue{Bound variables}: variables that are attached to $\lambda$ are {\em bound} variables 
--- $\lambda$ is a binder. 
\\
Bound variables in other systems: \vspace*{-0.5em}
\[
\begin{array}{c}
\lim_{x \to \infty} e^{-x} \\[0.3em]
\int_0^2 x^2 \, dx \\[0.3em]
\forall x, y, z: x \geq y \wedge y \geq z \Rightarrow x \geq z \\[0.3em]
\verb|int succ (x : int) { return x+1 ; }| 
\end{array}
\]
\end{itemize}
\end{slide}

\begin{slide}{Free and bound variables}
\begin{itemize}
\item \emblue{Free variables}: all those are not bound. We write $FV(e)$ for the set of free variables 
  of expression $e$.
\begin{wrap}
\vspace*{-0.5em}
\begin{eqnarray*}
FV(x) & = & \set{x} \\
FV(e \ e') & = & FV(e) \cup FV(e') \\
FV(\lambda x \ldot e) & = & FV(e) \setminus \set{x} \\[-2em]
\end{eqnarray*}
\end{wrap}
\Example (Free variables of $\lambda$-expressions) 
\[
\begin{array}{l}
FV(\lambda x \ldot x \, x) = \mbox{?} \\
FV((\lambda x \ldot x \, y) (\lambda y \ldot y \, z)) = \mbox{?} \\[-0.6em] 
\end{array}
\vspace*{-1em}
\]
\item A term without any free variable is a \emblue{closed} term; o.w. it's \emblue{open}.
\end{itemize}
\end{slide}

\begin{slide}{Variables and $\alpha$-equivalence}
\begin{itemize}
\item $\alpha$-equivalence: the names of bound variables are irrelevant --- they can be renamed.
\begin{wrap}
\vspace*{-0.8em}
\[
\lambda x \ldot e =_\alpha \lambda y \ldot e \set{y / x},
\mbox{ where } y \not\in FV(e)
\vspace*{-1em}
\]
\end{wrap}
\Example \vspace*{-1em}
\begin{eqnarray*}
\lambda x \ldot x & =_\alpha & \lambda y \ldot y \\ 
\lambda z \ldot x + z & \emred{\neq}_\alpha & \lambda z \ldot y + z \\
\lambda z \ldot x + z & \emred{\neq}_\alpha & \lambda x \ldot x + x \\[-4em]
\end{eqnarray*}
\item Formal definition of renaming: \vspace*{-0.5em}
\[
\begin{array}{l@{\;}l}
x \set{y / x} & \defeq y \\
z \set{y / x} & \defeq z,  \mbox{ if } z \neq x \\
(e \, e') \set{y / x} & \defeq (e \set{y / x}) (e' \set{y / x}) \\
(\lambda x \ldot e) \set{y / x} & \defeq \lambda x \ldot e \\
(\lambda z \ldot e) \set{y / x} & \defeq \lambda z \ldot (e \set{y / x})
\end{array}
\]
\end{itemize}
\end{slide}

\begin{slide}{$\beta$-reduction}
How do we compute in $\lambda$-calculus?
\begin{itemize}
\item In models like $\lambda$-calculus, computations are done by \emred{reductions}, 
  which basically turns an expression into another ``equivalent'' form. 
\item In $\lambda$-calculus, the essential reduction is called $\beta$-reduction:
\vspace*{-0.5em}
\begin{wrap}
\vspace*{-0.6em}
\emred{$$(\lambda x \ldot e) e' \reduce e [e' / x] \vspace*{-1.2em}$$}
\end{wrap}
\vspace*{-0.2em}
$e [e' / x]$ denotes the expression obtained from $e$ by \emblue{substituting} all free occurences 
of $x$ with $e'$. \vspace*{-0.5em}
\[
\begin{array}{l@{\;}l}
x [e / x] & \defeq e \\
z [e / x] & \defeq z,  \mbox{ if } z \neq x \\
(e_1 \, e_2) [e / x] & \defeq (e_1 [e / x]) (e_2 [e / x]) \\
(\lambda x \ldot e') [e / x] & \defeq \lambda x \ldot e' \\
(\lambda z \ldot e') [e / x] & \defeq \lambda z \ldot (e' [e / x]) \quad \mbox{if \emred{$z \not\in FV(e)$}} \\
(\lambda z \ldot e') [e / x] & \defeq \lambda y \ldot (e' \set{y / z} [e / x]) \quad \mbox{if \emred{$z \in FV(e)$} and $y $ is {\em fresh}} \\
\end{array}
\]
\end{itemize}
\end{slide}

\begin{slide}{$\beta$-reduction}
\begin{itemize}
\item A term of the form $(\lambda x \ldot e) e'$ is a \emblue{$\beta$-redex}. \\
$\beta$-reduction examples: \vspace*{-0.5em}
\[
\begin{array}{l}
(\lambda x \ldot x + 1) \, 2 \reduce 2 + 1 \\[0.1em]
(\lambda f . \lambda x \ldot f \, (f \, x)) \, \pmb{incr} 
\reduce \lambda x \ldot \pmb{incr} (\pmb{incr} \, x) \\[0.1em]
(\lambda f . \lambda g . \lambda x \ldot f \, (g \, x)) \, \pmb{incr} \, \pmb{double} \\[0.1em]
\qquad \reduce (\lambda g . \lambda x \ldot \pmb{incr} (g \, x)) \, \pmb{double} 
\reduce \lambda x \ldot \pmb{incr} (\pmb{double} \, x) \\[-1.8em]
\end{array}
\]
\item Single-step $\beta$-reduction can take place inside $\lambda$-expressions:
\vspace*{-0.5em}
\begin{wrap}
\vspace*{-0.5em}
\[
\begin{array}{c}
\begin{prooftree}
e_1 \reduce e_1'
\justifies 
e_1 e_2 \reduce e_1' e_2 
\end{prooftree}
\qquad 
\begin{prooftree}
e_2 \reduce e_2'
\justifies 
e_1 e_2 \reduce e_1 e_2'
\end{prooftree}
\qquad 
\begin{prooftree}
e \reduce e'
\justifies 
\lambda x \ldot e \reduce \lambda x \ldot e' 
\end{prooftree}
\end{array}
\vspace*{-0.5em}
\]
\end{wrap}
\Example \vspace*{-1em}
\begin{eqnarray*}
& & (\lambda f \ g \ x \ldot f (g \, x)) \, (\lambda x \ldot (\lambda  y \ldot x + y)  1)
  \, ((\lambda x \ y \ldot x * y) \, 2) 
\\
& \reduce & (\lambda f \ g \ x \ldot f (g \, x)) \, (\lambda x \ldot (\lambda y \ldot x + y)  1) 
  \, \empurple{(\lambda y \ldot 2 * y)}
\\
& \reduce & (\lambda f \ g \ x \ldot f (g \, x)) \, \empurple{(\lambda x \ldot x + 1)}
  \, (\lambda y \ldot 2 * y)
\\
& \reduce & \ldots \; \ldots 
\end{eqnarray*}
\end{itemize}
\end{slide}

\begin{slide}{$\beta$-reduction}
\begin{itemize}
\item Multi-step $\beta$-reduction $\reduceto_\beta$ is the {\em reflexive} and {\em transitive} closure 
  of $\betareduce$. 
\vspace*{-0.5em}
$$
(\lambda f \ g \ x \ldot f \ (g \ x)) \ \pmb{incr} \ \pmb{double} 
\reduceto_\beta \lambda x \ldot \pmb{incr} \ (\pmb{double} \ x)
$$ 
$\beta$-equivalence $=_\beta$ is the {\em symmetric}, reflexive and transitive closure. %%  of $\betareduce$.
\vspace*{-0.5em}
$$
(\lambda f \ g \ x \ldot f \ (g \ x)) \ \pmb{incr} \ \pmb{double} 
=_\beta (\lambda f \ g \ x \ldot g \ (f \ x)) \ \pmb{double} \ \pmb{incr} 
\vspace*{-2em}
$$ 
\item A term without any $\beta$-redexes is a \emblue{$\beta$-normal form}.
\vspace*{-0.5em}
$$ x \ y, \; \lambda x \ y \ldot y, \; y (\lambda x \ldot x) (\lambda x \ldot x) 
\vspace*{-2em}
$$
\item $\beta$-reductions of untyped $\lambda$-terms can be \emblue{infinite}.
\\[0.2em]
\Example : $\Omega \defeq (\lambda x \ldot x x) (\lambda x \ldot x x)$
\vspace*{-0.5em}
\[
(\lambda x \ldot x x) (\lambda x \ldot x x) 
\reduce (\lambda x \ldot x x) (\lambda x \ldot x x)
\reduce (\lambda x \ldot x x) (\lambda x \ldot x x) 
\reduce \ldots
\vspace*{-0.5em}
\]
The reduction can produce larger terms: \vspace*{-0.5em}
$$
(\lambda x \ldot x \ x \ x) (\lambda x \ldot x \ x \ x) 
\reduce (\lambda x \ldot x \ x \ x) (\lambda x \ldot x \ x \ x) (\lambda x \ldot x \ x \ x)
\reduce \cdots
$$ 
\end{itemize}
\end{slide}

\section{Programming $\lambda$-calculus}

\begin{slide}{Programming untyped $\lambda$-calculus}
In previous examples we have used $+$ and numerals ---  
  they are not contained in the pure $\lambda$-calculus, but can be represented by $\lambda$-terms.
\begin{itemize}
\item This is about \emblue{encoding}, just as we can represent data and computations in binary.
\item In terms of computability, \empurple{$\lambda$-calculus is equivalent to the Turing machine model}.
\begin{itemize}
\item More precisely, both $\lambda$-calculus and Turing machines define the same class of computable 
  (numeric) functions.
\item We shall see how different types of data and related operations can be programmed 
  in $\lambda$-calculus.
\end{itemize}
\end{itemize}
\end{slide}

\begin{slide}{Booleans}
\begin{itemize}
\item Normally, we use only \emblue{closed terms} to encode data and operations.
\begin{eqnarray*}
\pmb{True} & \defeq & \lambda x . \lambda y \ldot x  \\
\pmb{False} & \defeq & \lambda x . \lambda y \ldot y
\\[-3.5em]
\end{eqnarray*}
\item Boolean operations: 
\begin{eqnarray*}
\pmb{not} & \defeq & \lambda u \ldot \lambda v \ldot \lambda w \ldot u w v 
\\
\pmb{and} & \defeq & \lambda u \ldot \lambda v \ldot u v (\lambda x \ldot \lambda y \ldot y) 
\\
\ldots & \ldots &
\end{eqnarray*}
\vspace*{-1.5em} 
\[
\pmb{not} \ \pmb{True} = (\lambda u \ v \ w \ldot u \ w \ v) (\lambda x \ y \ldot x)
\reduce \lambda v \ w \ldot (\lambda x\ y \ldot x) w \ v 
\reduceto \lambda v \ w \ldot w
\vspace*{-2em}
\]
\item If-Then-Else: 
\[
\begin{array}{l}
\pmb{IF} \defeq \lambda u . \lambda v . \lambda w \ldot u v w
\\[0.5em]
\pmb{IF} \, \pmb{True} \, e_1 \, e_2 
= (\lambda u \ v \ w \ldot u v w) (\lambda x \ y \ldot x) e_1 e_2
\reduceto (\lambda x \ y \ldot x) e_1 e_2 \betareduce e_1
\end{array}
\]
\end{itemize}
\end{slide}

\begin{slide}{Natural numbers}
Natural numbers \vspace*{-0.5em}
\begin{eqnarray*}
\overline 0 & \defeq & \lambda f x \ldot x ,
\\
\overline 1 & \defeq & \lambda f x \ldot f x ,
\\ 
\overline 2 & \defeq & \lambda f x \ldot f (f x) ,
\\
\overline 3 & \defeq & \lambda f x \ldot f (f (f x)) ,
\\
& \ldots & \ldots
\end{eqnarray*}
\vspace*{-2.5em}
\begin{itemize}
\item This encoding is called the \emblue{Church numerals}: 
$$\emblue{\underline n \defeq \lambda f . \lambda x \ldot f^n x} \vspace*{-2em}$$
\item Natural number operations:
\begin{eqnarray*}
\pmb{add} & \defeq & \lambda u_1 . \lambda u_2 . \lambda f . \lambda x \ldot u_1 f (u_2 f x) 
\\
\pmb{mult} & \defeq & \lambda u_1 . \lambda u_2 . \lambda f . \lambda x \ldot u_1 (u_2 f) x 
\\
& \ldots & \ldots
\end{eqnarray*}
\end{itemize}
\end{slide}

\begin{slide}{Pairs and tuples}
\begin{itemize}
\item Pairs 
\[
\cpair{e_1, e_2} \defeq \lambda z \ldot z \ e_1 \ e_2
\]
Projections: \vspace*{-1.5em}
\begin{eqnarray*}
\pmb{first} & \defeq & \lambda u \ldot u (\lambda x \ y \ldot x) \\
\pmb{second} & \defeq & \lambda u \ldot u (\lambda x \ y \ldot y) \\[-4.5em]
\end{eqnarray*}
\item Tuples
\[
\cpair{e_1, \ldots, e_n} \defeq \lambda z \ldot z \ e_1 \cdots e_n
\]
$i$-th projection: 
$$\pmb{proj}_i \defeq \lambda u \ldot u (\lambda x_1 \cdots x_n \ldot x_i)$$
\end{itemize}
\end{slide}

\begin{slide}{Subtraction}
\begin{itemize}
\item We only concern unsigned arithmetic
 \[
\pmb{sub} \, n \, m = \left\{
\begin{array}{l@{\quad}l}
n - m & \mbox{if $n > m$} \\
0 & \mbox{otherwise}
\end{array}
\right .
\]
\emgreen{\bf Question:} what about signed arithmetic?
\item Define the predecessor: 
 \[
\pmb{pred} \, n = \left\{
\begin{array}{l@{\quad}l}
n - 1 & \mbox{if $n > 0$} \\
0 & \mbox{otherwise}
\end{array}
\right .
\]
Encoding in $\lambda$-calculus (using pair)
\[
\emblue{\pmb{pred} \defeq \lambda n \ldot \pmb{first} \ (n \ P \ \cpair{0, 0})}
\]
where $P \defeq \lambda x \ldot \cpair{\pmb{second} \ x, (\pmb{second} \ x) + 1}$. 
\end{itemize}
\end{slide}

\begin{slide}{Subtraction}
Programming \emblue{subtraction}
\begin{itemize}
\item Define $\pmb{sub}$ with $\pmb{pred}$:
\[
\emblue{\pmb{sub}\defeq \lambda n \ m \ldot m \ \pmb{pred} \ n}
\vspace*{-2em}
\]
\item An alternative definition of $\pmb{pred}$ (without using pair):
\[
\pmb{pred} \defeq \lambda n \ f \ x \ldot n \ T \ (\lambda u \ldot x) (\lambda v \ldot v) 
\]
where $T \defeq \lambda g \ h \ldot h \ (g \ f)$.
\\
It can be checked that for $n > 0$, 
\[
T^n (\lambda u \ldot x) = \lambda h \ldot h \ (f^{n-1} (x))
\]
\end{itemize}
\end{slide}

\begin{slide}{Inductive data structure --- list}
Programming lists
\begin{itemize}
\item List is the very basic data structure in functional programming.
\begin{itemize}
\item A list is just a sequence of elements, e.g., $[3, 1, 2, 4]$.
\item It has two constructs: \emblue{$\pmb{nil}$} --- the empty list, 
  and \emblue{$\pmb{cons}$} --- the operation that produces a list by adding an element to 
  the head of another list, e.g., 
\[
\pmb{cons} \ 2 \ [3, 1, 2, 4] = [2, 3, 1, 2, 4]
\]
\end{itemize}
\item $\pmb{nil}$ and $\pmb{cons}$ in $\lambda$-calculus:
\begin{eqnarray*}
\pmb{nil} & \defeq & \lambda x \ y \ldot y \\
\pmb{cons} & \defeq & \lambda h . \lambda t . \lambda x \ y \ldot x \ h \ t
\end{eqnarray*}
\emgreen{Exercise}: program $\pmb{isnil}$, $\pmb{length}$, \ldots
\end{itemize}
\end{slide}

\begin{slide}{Recursion and fix-points}
\begin{itemize}
\item Fix-point and recursion:
\\[0.2em]
Curry's (a.k.a. $\pmb{\mathrm{Y}}$-combinator): 
$\pmb{\mathrm{Y}} \defeq \lambda f \ldot (\lambda x \ldot f (x x)) (\lambda x \ldot f (x x))$
\vspace*{-0.5em}
\[
\pmb{\mathrm Y} F =_\beta F (\pmb{\mathrm Y} F)
\vspace*{-1em}
\]
Turing's fix-point: $\Theta \defeq (\lambda x \ldot \lambda f \ldot f (x x f)) 
(\lambda x \ldot \lambda f \ldot f (x x f))$
\\[0.2em]
There exist other fix-point combinators.
\end{itemize}
\end{slide}

\begin{slide}{The Church-Rosser theorem}
\begin{itemize}
\item \emblue{[Church \& Rosser, 1936]} If $e \reduceto e_1$ and $e \reduceto e_2$, then 
there exists $e'$ such that $e_1 \reduceto e'$ and $e_2 \reduceto e'$.
\[
\begin{array}{c@{\qquad\qquad}c}
\xymatrix@+=1pc{
  & e \ar@{->>}[dl] \ar@{->>}[dr] & \\
  e_1 \ar@{-->>}[dr] & & e_2 \ar@{-->>}[dl] \\
  & e' &
}
&
\xymatrix@+=1pc{
  & e \ar[dl] \ar[dr] & \\
  e_1 \ar@{-->}[dr] & & e_2 \ar@{-->}[dl] \\
  & e' &
} 
\\
(CR) & (Diamond)
\end{array}
\]
\item Corollary
\begin{itemize}
\item If $e_1 =_\beta e_2$ then there exists $e'$ with $e_1, e_2 \reduceto e$.
\\[0.3em]
\item If $e$ is in $\beta$-normal form and $e =_\beta e'$, then $e' \reduceto e$.
\end{itemize}
\end{itemize}
\end{slide}

\begin{slide}{Proof of the CR theorem}
Define \emred{a ``parallel reduction'' relation $\gg$}:
\[
\begin{array}{c}
\begin{prooftree}
\justifies 
e \gg e
\using{(Refl)}
\end{prooftree}
\qquad 
\begin{prooftree}
e \gg e'
\justifies 
\lambda x \ldot e \gg \lambda x \ldot e'
\using{(Abs)}
\end{prooftree}
\qquad 
\begin{prooftree}
e_1 \gg e_1' \quad e_2 \gg e_2'
\justifies 
e_1 e_2 \gg e_1' e_2'
\using{(App)}
\end{prooftree}
\\[1.5em]
\begin{prooftree}
e_1 \gg e_1' \quad e_2 \gg e_2'
\justifies 
(\lambda x \ldot e_1) e_2 \gg e_1' [ e_2' / x]
\using{(\parallel-\beta)}
\end{prooftree}
\end{array}
\]
The proof is supported by the following lemmas:
\begin{itemize}
\item[--] Lemma 1: \empurple{$\gg$ satisfies the Diamond-property}, i.e., if $e \gg e_1$ and $e \gg e_2$, 
then there exists $e'$ with $e_1 , e_2 \gg e'$. 
\item[--] Lemma 2: If $e \gg e'$ and $t \gg t'$, then $e [t / x] \gg e' [t' / x]$.
\item[--] Lemma 3: $\reduceto$ is the transitive closure of $\gg$.
\end{itemize}
\end{slide}

%% \begin{slide}{$\eta$-conversion}

%% \end{slide}

\end{document}
 

